\scanpage{109}
\mspage{54}

(C'est un foncteur \unsure{p.r.} au triple $\widetilde X\mathrel{\text{---}}X\mathrel{\text{---}}Y$ avec $X,Y$ $0$-connexes, $\widetilde X$ un rev. universel; plus \unsure{généralt}, si $\unsure{T|Z}\unsure{\mathrel{\text{---}}|\longrightarrow}X$ avec $\unsure{T|Z}$ $1$-connexe, on pose
\[
(11)\qquad \unsure{\mathcal E}_{X/Y}^{\unsure{T|Z}}=\operatorname{Ker}\bigl(\mathcal E_X^{\unsure{T|Z}}\unsure{\mathrel{\text{---}}|\longrightarrow}\mathcal E_Y^{\unsure{T|Z}}\bigr)
\]
on a un foncteur en \inserted{\unsure{le triple}} $(\unsure{T|Z}\unsure{\mathrel{\text{---}}|\longrightarrow}X\longrightarrow Y)$ variable. En particulier: $\mathcal E_{X/Y}^{\unsure{\xi|\zeta}}$, $\unsure{\xi|\zeta}$ un $p^t$ géom. de $X$, $\mathcal E_{X/Y}^{\Omega}$, $\mathcal E_{X/A}^{\widetilde X}$ (si $Y=\operatorname{Spec} A$), $\mathcal E_{B/A}^{\widetilde B}$ $\ldots$

\margin{\uline{NB} $\mathcal E_{X/Y}^T$, pour $T$ variable, $X/Y$ fixé, est défini à isom non unique \unsure{!} près, et se note $\mathcal E_{X/Y}$ \unsure{simplem\textsuperscript{t}.}}

Un cas intéressant est celui où on dispose d'une ``suite exacte d'homotopie universelle'' \deleted{de} (en dim $1$) pour $X\longrightarrow Y$, alors la donnée (pour $X=Y$ donné) de $T\unsure{\mathrel{\text{---}}|\longrightarrow}X$, (avec $T$ $1$-connexe) peut s'interpréter par la donnée \unsure{du} \unsure{composé}

\margin{Car $\mathcal E_X^T\longrightarrow\mathcal E_Y^T$ est épi.}
\margin{108}
